Tenim una fotocèl·lula en la qual el càtode és fet d'un material alcalí que només pot emetre electrons per efecte fotoelèctric si els fotons tenen una energia superior a $1,20\ \mathrm{eV}$. Enviem sobre el càtode un feix de fotons format per $10^{7}$ fotons/s d'una longitud d'ona de llum verda de $500\ \mathrm{nm}$.

\begin{apartat}{1,25}
Quina energia cinètica tindran els electrons arrencats del càtode per aquesta llum verda?
\end{apartat}

\begin{apartat}{1,25}
Si en lloc de $10^{7}$ fotons/s sobre el càtode hi enviem un feix 10 vegades més intens ($10^{8}$ fotons/s), quins canvis es produiran en l'emissió dels electrons? Justifiqueu la resposta.
\end{apartat}

\dades{%
  $1\ \mathrm{eV} = 1,602\times 10^{-19}\ \mathrm{J}$.\\
  $c = 3,00\times 10^{8}\ \mathrm{m\,s^{-1}}$.\\
  $h = 6,63\times 10^{-34}\ \mathrm{J\,s}$.}
